
\documentclass{article}
\usepackage{colortbl}
\usepackage{makecell}
\usepackage{multirow}
\usepackage{supertabular}

\begin{document}

\newcounter{utterance}

\centering \large Interaction Transcript for game `bbh', experiment `snarks', episode 139 with reflexion{-}Qwen3.5{-}9B.
\vspace{24pt}

{ \footnotesize  \setcounter{utterance}{1}
\setlength{\tabcolsep}{0pt}
\begin{supertabular}{c@{$\;$}|p{.15\linewidth}@{}p{.15\linewidth}p{.15\linewidth}p{.15\linewidth}p{.15\linewidth}p{.15\linewidth}}
    \# & \multicolumn{2}{c}{Player} && \multicolumn{2}{c}{Game Master} \\
    \hline

    \theutterance \stepcounter{utterance}  
    & & & \multicolumn{4}{p{0.6\linewidth}}{
        \cellcolor[rgb]{0.9,0.9,0.9}{
            \makecell[{{p{\linewidth}}}]{
                \texttt{\tiny{[P1$\langle$GM]}}
                \texttt{You have attempted this same task before and did not succeed. These are the plans you wrote after those attempts.} \\
\texttt{Trial \#1: New plan: I incorrectly identified option (A) as sarcastic because I focused on the surface{-}level irony of "Thank you" paired with distrust, missing the crucial context of the second clause. In option (A), the speaker explicitly states they will *not* trust an unsourced post, which is a logical, albeit cynical, stance, not a sarcastic reversal of meaning. True sarcasm in this context requires saying the opposite of what is meant; option (B) claims the speaker *will* trust an unsourced post over solid advice, which is an illogical, absurd statement intended to mock the idea that unsourced internet opinions hold more weight than expert advice. Next time, I will look for the option that presents an illogical or absurd preference as a genuine statement, as that indicates the speaker is being sarcastic by pretending to value the unsourced opinion. I will select option (B).} \\
\\ 
\\ 
\texttt{Determine which of two sentences is sarcastic.} \\
\\ 
\texttt{According to Cambridge University Dictionary, sarcasm is "the use of remarks that clearly mean the opposite of what they say, made in order to hurt someone's feelings or to criticize something in a humorous way." Sarcastic sentences often contain satirical or ironic utterances, hyperboles, ambivalent or witty remarks.} \\
\\ 
\texttt{Q: Which statement is sarcastic?} \\
\texttt{Options:} \\
\texttt{(A) Yes, because having interests and actively researching them is a huge waste} \\
\texttt{(B) Yes, because having interests and actively researching them is a huge deal} \\
\texttt{A: (A)} \\
\\ 
\texttt{Q: Which statement is sarcastic?} \\
\texttt{Options:} \\
\texttt{(A) No one is going to disagree with you on this. Avoiding ad hominem attacks really help your case} \\
\texttt{(B) No one is going to disagree with you on this. Ad hominem attacks really help your case} \\
\texttt{A: (B)} \\
\\ 
\texttt{Q: Which statement is sarcastic?} \\
\texttt{Options:} \\
\texttt{(A) Consistency in the league's punishments? What do you think this is supposed to be, politics?} \\
\texttt{(B) Consistency in the league's punishments? What do you think this is supposed to be, moral?} \\
\texttt{A: (A)} \\
\\ 
\texttt{Q: Which statement is sarcastic?} \\
\texttt{Options:} \\
\texttt{(A) Thank you stranger on the internet, I won't trust your three word post without any sources over someone offering solid advice claiming to actually be from the country in question} \\
\texttt{(B) Thank you stranger on the internet, I will trust your three word post without any sources over someone offering solid advice claiming to actually be from the country in question} \\
\texttt{A:} \\
            }
        }
    }
     \\ \\

    \theutterance \stepcounter{utterance}  
    & \multicolumn{4}{p{0.6\linewidth}}{
        \cellcolor[rgb]{0.561,0.737,0.561}{
            \makecell[{{p{\linewidth}}}]{
                \texttt{\tiny{[P1$\rangle$GM]}}
                \texttt{Based on the logic established in your previous trials and the definition of sarcasm provided, here is the analysis:} \\
\\ 
\texttt{**Analysis:**} \\
\texttt{*   **Option (A):** "I won't trust your three word post... over someone offering solid advice..."} \\
\texttt{    *   This statement expresses a logical, rational stance. It values sources and solid advice over anonymous, unsourced internet posts. While it might sound cynical, it is not saying the opposite of what is meant; it is stating a genuine preference for reliability. Therefore, it is **not** sarcastic.} \\
\texttt{*   **Option (B):** "I will trust your three word post... over someone offering solid advice..."} \\
\texttt{    *   This statement claims to value an unsourced, three{-}word post *more* than expert advice. In reality, trusting anonymous internet posts over verified experts is illogical and absurd. By stating this absurd preference as if it were a serious opinion, the speaker is mocking the idea that "internet wisdom" is superior to expertise. This is the opposite of what a rational person would mean, fitting the definition of sarcasm perfectly.} \\
\\ 
\texttt{**Conclusion:**} \\
\texttt{Option (B) is the sarcastic statement because it presents an illogical preference (valuing unsourced posts over solid advice) as a genuine statement to criticize the reliance on such posts.} \\
\\ 
\texttt{A: (B)} \\
            }
        }
    }
    & & \\ \\

    \theutterance \stepcounter{utterance}  
    & & & \multicolumn{2}{p{0.3\linewidth}}{
        \cellcolor[rgb]{0.9,0.9,0.9}{
            \makecell[{{p{\linewidth}}}]{
                \texttt{\tiny{[GM$|$GM]}}
                \texttt{Target: (B)} \\
            }
        }
    }
    & & \\ \\

    \theutterance \stepcounter{utterance}  
    & & & \multicolumn{2}{p{0.3\linewidth}}{
        \cellcolor[rgb]{0.9,0.9,0.9}{
            \makecell[{{p{\linewidth}}}]{
                \texttt{\tiny{[GM$|$GM]}}
                \texttt{game\_result = LOSE} \\
            }
        }
    }
    & & \\ \\

    \theutterance \stepcounter{utterance}  
    & & \multicolumn{4}{p{0.6\linewidth}}{
        \cellcolor[rgb]{0.95,0.95,0.95}{
            \makecell[{{p{\linewidth}}}]{
                \texttt{\tiny{[P1$\leftrightarrow$P]}}
                \texttt{Post{-}episode Reflexion (not sent to GM):} \\
\texttt{New plan: I failed because I misidentified the sarcastic sentence in the final trial by focusing on the cynicism of Option (A) rather than recognizing the absurdity of Option (B). Sarcasm here is defined by stating the opposite of a rational truth; trusting an unsourced, three{-}word post over solid expert advice is illogical, making Option (B) the clear sarcastic choice. Next time, I will immediately evaluate which option presents an illogical or absurd preference as a genuine statement. I will select the option that contradicts common sense to mock the subject, which is Option (B).} \\
            }
        }
    }
    & \\ \\

\end{supertabular}
}

\end{document}
